\section{Results and Data Analysis}

Following the experimental methodology established in Chapter 4, this chapter presents the empirical results of the performance evaluation. The data illustrates the latency, scalability, and structural resilience of the \textit{Thin Edge} architecture within the Plant Up! ecosystem.

\subsection{End-to-End Latency Benchmark}

To evaluate the real-time responsiveness of the system, an end-to-end latency benchmark was conducted to measure the interval between a physical sensor trigger and its final persistence in the cloud database. As the MQTT protocol operates asynchronously, these measurements were derived using the differential timestamping method described in Section 4.3.

The results, as summarized in Table \ref{tab:end_to_end_latency}, indicate that while the cloud-centric approach introduces multiple network hops, the total response time remains well within the acceptable boundaries for botanical monitoring.

\begin{table}[htbp]
    \centering
    \begin{tabular}{@{}ll@{}}
        \toprule
        \textbf{Component} & \textbf{Measured Latency} \\ \midrule
        ESP32 Measurement to MQTT Publish & 24\,ms \\
        MQTT Agent to Supabase Ingress & 142\,ms \\
        Supabase Real-time to Database & 98\,ms \\ \midrule
        \textbf{Total End-to-End Latency} & \textbf{264\,ms} \\ \bottomrule
    \end{tabular}
    \caption{End-to-End Ingestion Latency (Sensors to Persistence). Source: Own compilation.}
    \label{tab:end_to_end_latency}
\end{table}

The observed latency of approximately 264\,ms reflects the impact of the transport layer and the centralized gateway logic. While this duration is significantly higher than localized edge processing ($<$50\,ms), it provides the foundational infrastructure required for the resource-heavy AI and social features that define the Plant Up! experience.

\subsection{Database Performance: TimescaleDB vs. Postgres}

A critical component of backend performance is the efficiency of telemetry retrieval. To quantify the benefit of the TimescaleDB hypertable implementation, a comparative query execution analysis was performed. Using the PostgreSQL \texttt{EXPLAIN ANALYZE} utility via the Supabase REST API, a standard relational table scan was compared against a chunked hypertable scan across a dataset comprising 500,000 historical sensor data points. The results represent an average computed across 100 consecutive benchmark iterations to ensure analytical validity.

\begin{table}[htbp]
    \centering
    \begin{tabular}{@{}ll@{}}
        \toprule
        \textbf{Query Type} & \textbf{Execution Time (Average)} \\ \midrule
        Standard Table Scan (B-Tree Index) & 74.8\,ms \\
        TimescaleDB Hypertable (Chunking) & 8.6\,ms \\ \midrule
        \textbf{Performance Improvement} & \textbf{8.7x faster (870\,\%)} \\ \bottomrule
    \end{tabular}
    \caption{Query Performance Comparison for Time-Series Data (500k rows, 100 runs). Source: Own compilation.}
    \label{tab:timescaledb_performance}
\end{table}

The results in Table \ref{tab:timescaledb_performance} confirm that hypertable chunking significantly reduces the I/O burden during range-based queries. By limiting the search space to active temporal chunks, the system maintains high retrieval speeds even as the total database volume scales, preventing the performance degradation typically associated with traditional relational architectures. At an 8.7x performance multiplier, the TimescaleDB engine proves highly effective in optimizing temporal metric retrieval.

While hypertable chunking ensures rapid retrieval of stored data (reads), creating actionable insights requires stable initial data ingress (writes). Therefore, the ingestion framework's reliability needed to be empirically validated.

\subsection{Throughput and Scalability Stress Test}

To evaluate the scalability of the middleware gateway pattern, a stress test was executed simulating a high-velocity telemetry influx. Using a Python-based simulation environment, a fleet of 50 concurrent devices were orchestrated to rapidly publish 10,000 payloads to the MQTT broker. This scenario simulates a ``Disaster Recovery'' spike, representing multiple edge nodes flushing local caches simultaneously after regaining network connectivity.

Under a load of 10,000 consecutive writes, the system successfully persisted all payloads with a 0\% drop rate. This benchmark proves that the MQTT gateway's concurrent queuing and batch processing logic effectively limits the frequency of outbound HTTP requests. By pooling the insertion requests, the system prevents hitting the cloud provider's API rate limits and significantly reduces network connection handshake overhead.

\newpage
\subsection{Architectural Comparison Summary}

The overarching premise of this evaluation is to quantify the necessary trade-offs introduced by adopting a \textit{Thin Edge} topology. Table \ref{tab:architecture_comparison} contextualizes the measured Plant Up! benchmarks against the recognized theoretical baselines of traditional, localized architectures.

\begin{table}[htbp]
    \centering
    \renewcommand{\arraystretch}{1.4}
    \begin{tabular}{@{} l l l @{}}
        \toprule
        \textbf{Metric} & \textbf{Traditional Edge} & \textbf{Plant Up! (Thin Edge)} \\ \midrule
        \textbf{Response Flexibility} & Low (Firmware Bound) & High (Cloud Deployed) \\
        \textbf{Analysis Power} & Low (ESP32 Limits) & Highly Scalable (Cloud Server) \\
        \textbf{Connectivity Dependency} & Low (Autonomous) & High (Network Required) \\
        \textbf{Latency Benchmark} & $\approx$10\,ms (Theoretical) & 264\,ms (Empirical) \\ \bottomrule
    \end{tabular}
    \caption{Contextual Analysis of Measured Thin Edge Performance vs. Localized Baselines. Source: Own compilation.}
    \label{tab:architecture_comparison}
\end{table}

While localized systems offer superior theoretical raw latency and strict offline autonomy, the empirical data gathered confirms that the Plant Up! infrastructure successfully mitigates its inherent networking delays. By leveraging database temporal chunking and ingestion batching, the platform provides a highly robust foundation for advanced botanical insights, justifying the trade-off of sub-second reaction times in exchange for scalable, data-driven operations.

